\documentclass[11pt,a4paper]{article}

% =====================================================================
% Codificación e idioma
% =====================================================================
\usepackage[utf8]{inputenc}
\usepackage[T1]{fontenc}
\usepackage[spanish,es-noshorthands]{babel}
\usepackage{lmodern}
\usepackage{textcomp}

% =====================================================================
% Layout
% =====================================================================
\usepackage[margin=2.5cm]{geometry}
\usepackage{parskip}
\setlength{\parindent}{0pt}
\setlength{\parskip}{0.55em plus 0.2em minus 0.1em}

% =====================================================================
% Tablas, listas, color, gráficos
% =====================================================================
\usepackage{booktabs}
\usepackage{array}
\usepackage{enumitem}
\usepackage{xcolor}
\usepackage{graphicx}
\usepackage{tikz}
\usetikzlibrary{positioning, arrows.meta, shapes.geometric, calc, fit, backgrounds}
\usepackage{caption}
\usepackage{subcaption}
\usepackage{tcolorbox}
\tcbuselibrary{skins, breakable}

% --- Colores para cajas resaltadas ---
\definecolor{noteborder}{rgb}{0.30,0.45,0.70}
\definecolor{notebg}{rgb}{0.94,0.96,1.00}
\definecolor{warnborder}{rgb}{0.75,0.45,0.10}
\definecolor{warnbg}{rgb}{1.00,0.97,0.90}

\newtcolorbox{idea}[1][]{
    enhanced, breakable,
    colback=notebg, colframe=noteborder,
    boxrule=0.5pt, arc=2pt,
    left=8pt, right=8pt, top=4pt, bottom=4pt,
    title={\bfseries #1}, fonttitle=\bfseries,
    coltitle=noteborder, colbacktitle=notebg,
    boxed title style={frame hidden, colback=notebg}
}

\newtcolorbox{ojo}[1][]{
    enhanced, breakable,
    colback=warnbg, colframe=warnborder,
    boxrule=0.5pt, arc=2pt,
    left=8pt, right=8pt, top=4pt, bottom=4pt,
    title={\bfseries #1}, fonttitle=\bfseries,
    coltitle=warnborder, colbacktitle=warnbg,
    boxed title style={frame hidden, colback=warnbg}
}

% =====================================================================
% Listings de código
% =====================================================================
\usepackage{listings}
\definecolor{codegray}{rgb}{0.96,0.96,0.96}
\definecolor{codeborder}{rgb}{0.85,0.85,0.85}
\definecolor{codecomment}{rgb}{0.30,0.55,0.30}
\definecolor{codekeyword}{rgb}{0.10,0.30,0.70}
\definecolor{codestring}{rgb}{0.65,0.15,0.10}

\lstdefinestyle{shellstyle}{
    language=bash,
    backgroundcolor=\color{codegray},
    basicstyle=\ttfamily\small,
    breaklines=true,
    breakatwhitespace=false,
    frame=single,
    rulecolor=\color{codeborder},
    showstringspaces=false,
    keywordstyle=\color{codekeyword}\bfseries,
    commentstyle=\color{codecomment}\itshape,
    stringstyle=\color{codestring},
    aboveskip=1em,
    belowskip=1em,
    columns=fullflexible,
    upquote=true,
}

\lstdefinestyle{cstyle}{
    language=C,
    backgroundcolor=\color{codegray},
    basicstyle=\ttfamily\footnotesize,
    breaklines=true,
    frame=single,
    rulecolor=\color{codeborder},
    showstringspaces=false,
    keywordstyle=\color{codekeyword}\bfseries,
    commentstyle=\color{codecomment}\itshape,
    stringstyle=\color{codestring},
    numbers=left,
    numberstyle=\tiny\color{gray},
    numbersep=8pt,
    aboveskip=1em,
    belowskip=1em,
    columns=fullflexible,
    upquote=true,
}

\lstdefinestyle{asmstyle}{
    backgroundcolor=\color{codegray},
    basicstyle=\ttfamily\footnotesize,
    breaklines=true,
    frame=single,
    rulecolor=\color{codeborder},
    showstringspaces=false,
    commentstyle=\color{codecomment}\itshape,
    aboveskip=1em,
    belowskip=1em,
    columns=fullflexible,
}

\lstdefinestyle{outstyle}{
    backgroundcolor=\color{white},
    basicstyle=\ttfamily\small,
    frame=single,
    rulecolor=\color{codeborder},
    breaklines=true,
    aboveskip=1em,
    belowskip=1em,
    columns=fullflexible,
}

\lstset{style=shellstyle}

% =====================================================================
% Hyperref
% =====================================================================
\usepackage[colorlinks=true,
            linkcolor=blue!60!black,
            urlcolor=blue!50!black,
            citecolor=blue!60!black]{hyperref}

% =====================================================================
% Portada
% =====================================================================
\title{
    \vspace{-2.5em}
    \textbf{Resolución del Práctico 2} \\[0.3em]
    \large Compilación, linking y layout de memoria de un proceso
}
\author{
    Tomás Rodeghiero \\
    \small Sistemas Operativos --- Universidad Nacional de Río Cuarto
}
\date{2026}

\begin{document}

\maketitle

\begin{abstract}
\noindent
Este documento es la resolución comentada del Práctico 2 de Sistemas
Operativos (UNRC, 2026), centrado en el \emph{toolchain} de C: las
etapas de pre-procesado, compilación, ensamblado y linking, junto con
la diferencia entre \textbf{bibliotecas estáticas y dinámicas} y el
\textbf{layout en memoria} de un proceso (text, data, bss, heap,
stack). El objetivo es que un estudiante pueda leer este PDF de
corrido y entender qué hace cada herramienta y por qué hace falta.
Los comandos fueron ejecutados sobre macOS arm64 (formato Mach-O); el
razonamiento es portable: en GNU/Linux las mismas ideas se manifiestan
sobre ELF, con \texttt{readelf}/\texttt{ldd}/\texttt{objdump} donde acá
usamos \texttt{otool}/\texttt{nm}/\texttt{objdump}.
\end{abstract}

\tableofcontents
\bigskip

% =====================================================================
\section{Introducción: el toolchain y la memoria de un proceso}
% =====================================================================

Programar en C parece resumirse a un único comando: \texttt{gcc -o
myprog main.c hello.c}. Detrás de esa línea se ejecuta una secuencia
ordenada de programas que se conoce como \emph{toolchain}, y el
resultado final ---el ejecutable--- se organiza en memoria con un
layout muy concreto. Antes de los ejercicios conviene fijar el
vocabulario, porque las preguntas del práctico se entienden mucho
mejor si se ve el panorama completo.

\subsection*{Las cuatro etapas del toolchain}

Las Notas del curso descomponen la compilación en cuatro etapas:

\begin{enumerate}[leftmargin=*]
    \item \textbf{Pre-procesado (\texttt{cpp}).} Expande las
        directivas \texttt{\#include}, \texttt{\#define} y demás
        macros. Salida: \texttt{.i} (C ya expandido, todavía texto).
        Se invoca con \texttt{gcc -E}.
    \item \textbf{Compilación (\texttt{cc1}).} Traduce el C
        pre-procesado a \emph{assembly} para la arquitectura objetivo.
        Salida: \texttt{.s}. Se invoca con \texttt{gcc -S}.
    \item \textbf{Ensamblado (\texttt{as}).} Convierte el assembly en
        un \emph{archivo objeto} binario (\texttt{.o}) en el formato
        del SO: \textbf{ELF} en GNU/Linux, \textbf{Mach-O} en macOS,
        \textbf{COFF/PE} en Windows. Se invoca con \texttt{gcc -c}.
    \item \textbf{Linking (\texttt{ld}).} Combina archivos objeto y
        bibliotecas, recalcula direcciones y resuelve referencias
        externas para producir el ejecutable final.
\end{enumerate}

Las opciones \texttt{-E}, \texttt{-S} y \texttt{-c} detienen el
proceso después de la primera, segunda y tercera etapa
respectivamente.

\subsection*{¿Qué hay dentro de un archivo objeto?}

Un \texttt{.o} no es código binario plano: es un contenedor con varias
\textbf{secciones}:

\begin{itemize}[leftmargin=*]
    \item \texttt{.text}: código máquina ejecutable.
    \item \texttt{.data}: variables globales/estáticas con valor
        inicial distinto de cero.
    \item \texttt{.bss}: variables globales/estáticas inicializadas a
        cero (no ocupan lugar en disco, solo se reservan los bytes en
        memoria).
    \item \texttt{.rodata}: constantes y literales de solo lectura.
        En Mach-O, los strings caen en la sección
        \texttt{\_\_TEXT,\_\_cstring}.
    \item \textbf{Tabla de símbolos:} lista las funciones y variables
        definidas o referenciadas por el módulo. Cada símbolo es
        \textit{local}, \textit{global definido} o \texttt{*UND*} (no
        resuelto).
    \item \textbf{Tabla de reubicación:} indica qué instrucciones
        tienen operandos pendientes de completar y a qué símbolo se
        refieren. Es el mapeo $\texttt{reloc(addr, tipo)} \to
        \texttt{símbolo}$ que mencionan las notas.
\end{itemize}

\subsection*{¿Qué hace exactamente el linker?}

El linker realiza dos tareas conceptuales:

\begin{enumerate}[leftmargin=*]
    \item Concatenar las secciones de los objetos
        (\texttt{.text}+\texttt{.text}+\ldots) y \emph{recalcular} las
        direcciones absolutas.
    \item Resolver las referencias externas, completando las
        instrucciones que la tabla de reubicación marcaba como
        pendientes.
\end{enumerate}

Las llamadas a funciones definidas en el mismo ejecutable se resuelven
en \textbf{linking estático}; las que viven en una biblioteca
compartida (como \texttt{printf} en la libc) se resuelven en
\textbf{linking dinámico}, en tiempo de \emph{carga} o de
\emph{primera llamada}, mediante \emph{stubs} y la \emph{Global Offset
Table} (GOT).

\subsection*{Layout en memoria de un proceso}

Una vez que el ejecutable se carga en memoria, el espacio de
direcciones del proceso queda organizado en regiones bien definidas.
Las secciones del archivo objeto se mapean (con permisos de lectura,
escritura, ejecución) a estas regiones del proceso:

\begin{center}
\begin{tikzpicture}[font=\ttfamily\small,
    region/.style={draw, minimum width=5cm, minimum height=0.85cm,
                   align=center, font=\ttfamily\small}]
    \node[region, fill=blue!10]                   (text) at (0, 0)    {\textbf{text} \\ (\texttt{.text}, código)};
    \node[region, fill=blue!10, below=0.05cm of text]    (rodata) {\textbf{rodata} \\ (\texttt{.rodata} / cstring)};
    \node[region, fill=green!12, below=0.05cm of rodata] (data) {\textbf{data} \\ (\texttt{.data}, globales con valor)};
    \node[region, fill=green!12, below=0.05cm of data]   (bss)  {\textbf{bss} \\ (\texttt{.bss}, globales en cero)};
    \node[region, fill=yellow!25, below=0.4cm of bss,
          minimum height=1.2cm]                          (heap) {\textbf{heap} \\ (\texttt{malloc}, crece $\uparrow$)};
    \node[region, fill=red!12,  below=1.7cm of heap,
          minimum height=1.4cm]                          (stack) {\textbf{stack} \\ (variables locales, crece $\downarrow$)};

    % flechas y direcciones
    \draw[->, thick] ($(heap.east)+(0.2,-0.3)$) -- ++(0,0.6) node[right, font=\small\sffamily]{crece};
    \draw[->, thick] ($(stack.east)+(0.2,0.3)$) -- ++(0,-0.6) node[right, font=\small\sffamily]{crece};

    \node[anchor=west, font=\small\sffamily] at ($(text.west)+(-1.6,0)$)  {dir.\ bajas};
    \node[anchor=west, font=\small\sffamily] at ($(stack.west)+(-1.6,0)$) {dir.\ altas};
\end{tikzpicture}
\end{center}

\begin{idea}[Cómo mapea cada sección al proceso]
\begin{itemize}[leftmargin=*, itemsep=0pt]
    \item \texttt{.text} y \texttt{.rodata} se mapean como
        \textbf{r-x} y \textbf{r{-}{-}} (no escribibles).
    \item \texttt{.data} y \texttt{.bss} como \textbf{rw-}.
    \item El \textbf{stack} crece hacia direcciones \emph{bajas}: cada
        \texttt{call} empuja un nuevo \emph{stack frame} hacia abajo.
    \item El \textbf{heap} crece hacia direcciones \emph{altas}: cada
        \texttt{malloc()} pide al kernel más espacio vía \texttt{brk}
        o \texttt{mmap}.
    \item Stack y heap crecen uno hacia el otro pero nunca se tocan:
        el SO los aísla con páginas de guarda.
\end{itemize}
\end{idea}

\subsection*{Bibliotecas estáticas vs dinámicas}

Es la dicotomía central del práctico:

\begin{itemize}[leftmargin=*]
    \item Una biblioteca \textbf{estática} (\texttt{.a}) es un archivo
        \texttt{ar} con varios \texttt{.o} y un índice de símbolos.
        El linker copia el código necesario \emph{dentro} del
        ejecutable. Resultado: un binario más grande pero
        autocontenido.
    \item Una biblioteca \textbf{dinámica}
        (\texttt{.so}/\texttt{.dylib}/\texttt{.dll}) vive en un
        archivo aparte. El ejecutable solo guarda una referencia y el
        \emph{loader} del SO la mapea en el proceso al arrancar (o
        bajo demanda con \texttt{dlopen}). Si varios procesos la usan,
        \emph{una sola copia} en memoria los sirve a todos.
\end{itemize}

Una tabla resumida (los detalles aparecen en los ejercicios 7 y 8):

\begin{center}
\small
\begin{tabular}{p{4.0cm} p{4.8cm} p{4.8cm}}
\toprule
 & \textbf{Estática (\texttt{.a})} & \textbf{Dinámica (\texttt{.so}/\texttt{.dylib})} \\
\midrule
Cuándo se resuelven los símbolos & En tiempo de \emph{linking} & En tiempo de \emph{carga} o primera llamada \\
Qué queda en el ejecutable & Una copia del código de la biblioteca & Solo una referencia (stub + GOT) \\
Tamaño del binario & Más grande & Más chico \\
Actualizar la biblioteca & Hay que relinkear & Basta con reemplazar el archivo \\
Memoria entre procesos & Una copia por proceso & Una sola copia compartida \\
Sintaxis típica & \texttt{ar rcs libfoo.a a.o b.o} & \texttt{gcc -fPIC -shared} \\
\bottomrule
\end{tabular}
\end{center}

Con esto en la cabeza, los ejercicios del práctico se vuelven una
visita guiada por cada uno de estos componentes.

% =====================================================================
\section{Ejercicio 1 --- Pre-procesado, assembly y objeto de \texttt{hello.c}}
% =====================================================================

\textbf{Enunciado.} Tomar el archivo \texttt{hello.c}, generar las
salidas intermedias del compilador (pre-procesado, assembly y archivo
objeto) y entender qué información contiene cada una.

\textbf{Programa fuente:}

\begin{lstlisting}[style=cstyle]
/* hello.c */
#define hi "Hello world"

char *hello(void)
{
    return hi;
}
\end{lstlisting}

\textbf{Análisis.} \texttt{hi} no es una variable: es una
\textbf{macro}. El preprocesador la reemplaza textualmente por la
cadena \texttt{"Hello world"} antes de que el compilador siquiera vea
el código. Por eso \texttt{hi} no va a aparecer en ningún archivo
posterior al \texttt{.i}.

\subsection*{1.a Pre-procesado (\texttt{hello.i})}

\textbf{Procedimiento:}

\begin{lstlisting}
gcc -E hello.c -o hello.i
\end{lstlisting}

\textbf{Resultado.} El contenido relevante de \texttt{hello.i} es:

\begin{lstlisting}[style=cstyle]
# 1 "hello.c"
# 1 "<built-in>" 1
...
char *hello(void)
{
    return "Hello world";
}
\end{lstlisting}

\textbf{Conclusión.} El identificador \texttt{hi} desapareció: la
macro fue expandida y reemplazada por el literal \texttt{"Hello
world"}. Las líneas \texttt{\# n "archivo"} son marcadores que el
preprocesador deja para que los mensajes de error apunten al archivo
y línea originales.

\subsection*{1.b Assembly (\texttt{hello.s})}

\textbf{Procedimiento:}

\begin{lstlisting}
gcc -S hello.c -o hello.s
\end{lstlisting}

\textbf{Resultado:}

\begin{lstlisting}[style=asmstyle]
        .section        __TEXT,__text,regular,pure_instructions
        .globl  _hello
        .p2align        2
_hello:
        adrp    x0, l_.str@PAGE
        add     x0, x0, l_.str@PAGEOFF
        ret
        .section        __TEXT,__cstring,cstring_literals
l_.str:
        .asciz  "Hello world"
\end{lstlisting}

\textbf{Cómo leerlo, línea por línea:}
\begin{itemize}[leftmargin=*]
    \item \texttt{\_\_TEXT,\_\_text} contiene el código de
        \texttt{hello}. La etiqueta \texttt{\_hello} (con underscore
        delante, convención macOS) es global: la verá el linker.
    \item \texttt{\_\_TEXT,\_\_cstring} contiene el literal,
        identificado por la etiqueta local \texttt{l\_.str}. En ELF
        esto sería la sección \texttt{.rodata}.
    \item ARM64 codifica direcciones en páginas de 4\,KiB, por eso la
        dirección del literal se arma en dos pasos: \texttt{adrp}
        carga la dirección base de la página, y \texttt{add} suma el
        offset dentro de ella.
    \item El valor de retorno se deposita en \texttt{x0} (convención
        AAPCS64) y la función vuelve con \texttt{ret}, que salta a la
        dirección guardada en el \emph{Link Register} \texttt{x30}.
\end{itemize}

\subsection*{1.c Archivo objeto (\texttt{hello.o})}

\textbf{Procedimiento:}

\begin{lstlisting}
gcc -c hello.c -o hello.o
file hello.o
\end{lstlisting}

\textbf{Resultado:}

\begin{lstlisting}[style=outstyle]
hello.o: Mach-O 64-bit object arm64
\end{lstlisting}

\textbf{Conclusión.} \texttt{hello.o} es un binario \emph{relocatable}
en formato Mach-O (Linux usaría ELF). Ya contiene el código máquina de
\texttt{hello()}, pero \emph{su dirección sigue siendo \texttt{0x0}}:
los offsets quedan relativos al inicio de cada sección, y será el
linker el que les asigne direcciones absolutas al combinarlo con
otros objetos.

% =====================================================================
\section{Ejercicio 2 --- Compilar \texttt{main.c} a \texttt{main.o}}
% =====================================================================

\textbf{Enunciado.} Compilar el programa principal hasta el archivo
objeto (sin linkear) y analizar qué quedó adentro y qué no.

\begin{lstlisting}[style=cstyle]
/* main.c */
#include <stdio.h>            /* prototipo de printf */

extern char* hello(void);     /* hello vive en otro modulo */

int main(void)
{
    printf("%s\n", hello());
    return 0;
}
\end{lstlisting}

\textbf{Análisis.} Tres elementos importantes:

\begin{itemize}[leftmargin=*]
    \item \texttt{\#include <stdio.h>} introduce el \emph{prototipo}
        de \texttt{printf} para que el compilador valide los tipos de
        los argumentos. La \emph{definición} (el código de
        \texttt{printf}) vive en la libc del sistema.
    \item \texttt{extern char* hello(void)} es solo una
        \emph{declaración}: le dice al compilador ``existe esta
        función en algún otro \texttt{.o}, ya la encontrará el linker''.
    \item Como consecuencia, \texttt{printf} y \texttt{hello} van a
        aparecer en \texttt{main.o} como símbolos externos no
        resueltos (\texttt{*UND*}).
\end{itemize}

\textbf{Procedimiento:}

\begin{lstlisting}
gcc -c main.c -o main.o
file main.o
\end{lstlisting}

\textbf{Resultado:}

\begin{lstlisting}[style=outstyle]
main.o: Mach-O 64-bit object arm64
\end{lstlisting}

\textbf{Conclusión.} La opción \texttt{-c} detiene el proceso tras el
ensamblado: preprocesa, compila y ensambla, pero \emph{no} linkea.
\texttt{main.o} contiene el código de \texttt{main} y el literal
\texttt{"\%s\textbackslash n"}, pero las llamadas a \texttt{hello} y
\texttt{printf} están todavía sin resolver.

\begin{idea}[Por qué compilar a objetos]
Compilar a \texttt{.o} en lugar de directamente al ejecutable habilita
la \textbf{compilación incremental}: si más adelante se modifica
\texttt{hello.c}, basta con regenerar \texttt{hello.o} y volver a
linkear, sin tocar \texttt{main.c}. Esta es la propiedad que aprovecha
\texttt{make} y la base de cualquier sistema de \emph{build}
moderno.
\end{idea}

% =====================================================================
\section{Ejercicio 3 --- Análisis de \texttt{main.o} y \texttt{hello.o}}
% =====================================================================

\textbf{Enunciado.} Inspeccionar a fondo los dos archivos objeto:
secciones, tabla de símbolos, tabla de reubicación y desensamblado.
La idea es \emph{ver} el resultado de las etapas anteriores y entender
por qué hace falta un linker.

\subsection*{3.a Tipo de archivo}

\textbf{Procedimiento:}

\begin{lstlisting}
file main.o hello.o
\end{lstlisting}

\textbf{Resultado:}

\begin{lstlisting}[style=outstyle]
main.o:  Mach-O 64-bit object arm64
hello.o: Mach-O 64-bit object arm64
\end{lstlisting}

\textbf{Conclusión.} Ambos son objetos relocatables Mach-O para arm64.
En GNU/Linux la salida sería ``ELF 64-bit LSB relocatable, ARM
aarch64''. Las direcciones todavía no están resueltas.

\subsection*{3.b Secciones (\texttt{objdump -h}) y tamaños}

\textbf{Procedimiento:}

\begin{lstlisting}
objdump -h main.o hello.o
\end{lstlisting}

\textbf{Resultado:}

\begin{lstlisting}[style=asmstyle]
main.o:	file format mach-o arm64
Sections:
Idx Name             Size     VMA              Type
  0 __text           00000040 0000000000000000 TEXT
  1 __cstring        00000004 0000000000000040 DATA
  2 __compact_unwind 00000020 0000000000000048 DATA

hello.o:	file format mach-o arm64
Sections:
Idx Name             Size     VMA              Type
  0 __text           0000000c 0000000000000000 TEXT
  1 __cstring        0000000c 000000000000000c DATA
  2 __compact_unwind 00000020 0000000000000018 DATA
\end{lstlisting}

\textbf{Análisis.}

\begin{center}
\small
\begin{tabular}{l c c c}
\toprule
\textbf{Objeto} & \texttt{.text} & \texttt{.data} & \texttt{\_\_cstring} \\
\midrule
\texttt{main.o}  & \texttt{0x40} (64 B) & --- (sección ausente) & \texttt{0x04}: \texttt{"\%s\textbackslash n"} \\
\texttt{hello.o} & \texttt{0x0c} (12 B) & --- (sección ausente) & \texttt{0x0c}: \texttt{"Hello world"} \\
\bottomrule
\end{tabular}
\end{center}

\texttt{main.o} es bastante más grande que \texttt{hello.o} en su
sección de código porque tiene prólogo y epílogo de \texttt{main},
armado de argumentos de \texttt{printf} y dos llamadas a funciones
(\texttt{hello}, \texttt{printf}). \texttt{hello()}, en cambio, solo
arma un puntero y retorna: tres instrucciones, 12 bytes.

\textbf{Conclusión.} No aparece sección \texttt{\_\_DATA,\_\_data} en
ninguno de los dos objetos: el programa \textbf{no tiene variables
globales con valor inicial distinto de cero}. Tampoco hay
\texttt{.bss} porque no hay globales inicializadas en cero. El literal
\texttt{"Hello world"} no va en \texttt{.data} sino en
\texttt{\_\_cstring} porque es de solo lectura.

\subsection*{3.c Contenido de cada sección}

\textbf{Procedimiento:}

\begin{lstlisting}
objdump -s main.o hello.o
\end{lstlisting}

\textbf{Resultado (parte relevante):}

\begin{lstlisting}[style=asmstyle]
main.o:
Contents of section __TEXT,__text:
 0000 ff8300d1 fd7b01a9 fd430091 08008052
 0010 e80b00b9 bfc31fb8 00000094 e8030091
 0020 000100f9 00000090 00000091 00000094
 0030 e00b40b9 fd7b41a9 ff830091 c0035fd6
Contents of section __TEXT,__cstring:
 0040 25730a00                             %s..

hello.o:
Contents of section __TEXT,__text:
 0000 00000090 00000091 c0035fd6           .........._.
Contents of section __TEXT,__cstring:
 000c 48656c6c 6f20776f 726c6400           Hello world.
\end{lstlisting}

\textbf{Lectura.} Los bytes \texttt{25 73 0a 00} son la secuencia
ASCII de \texttt{"\%s\textbackslash n\textbackslash 0"}: 0x25 = `\%',
0x73 = `s', 0x0a = \texttt{LF}, 0x00 = terminador. Los bytes en
\texttt{hello.o} forman directamente \texttt{"Hello world\textbackslash 0"}.

\subsection*{3.d Direcciones de \texttt{main()} y \texttt{hello()}}

\textbf{Procedimiento:}

\begin{lstlisting}
objdump --syms main.o hello.o
\end{lstlisting}

\textbf{Resultado (parte global):}

\begin{lstlisting}[style=asmstyle]
main.o:
 0000000000000000 g     F __TEXT,__text    _main

hello.o:
 0000000000000000 g     F __TEXT,__text    _hello
\end{lstlisting}

\textbf{Conclusión.} Ambas funciones empiezan en el offset
\texttt{0x0} de su sección \texttt{\_\_text}. Esto es lo
\emph{normal} en archivos objeto relocatables: cada módulo se piensa
como si viviera solo, en su propio espacio que arranca en cero. El
linker, más adelante, apila las secciones y reasigna direcciones
absolutas únicas.

\subsection*{3.e Dónde se almacena el literal apuntado por \texttt{hi}}

\textbf{Recordatorio.} \texttt{hi} es la macro \texttt{\#define hi
"Hello world"}: \emph{no} es una variable, así que no aparece como
símbolo en la tabla. Lo que sí aparece es el literal al que la macro
se expande:

\begin{lstlisting}[style=asmstyle]
000000000000000c l     O __TEXT,__cstring    l_.str
\end{lstlisting}

\textbf{Conclusión.} El string vive en la sección \texttt{\_\_cstring}
de \texttt{hello.o} (de solo lectura, equivalente a \texttt{.rodata}
en ELF), con etiqueta interna \texttt{l\_.str}, en el offset
\texttt{0x0c}. La función \texttt{hello()} arma esa dirección con
\texttt{adrp+add} y la devuelve en \texttt{x0}.

\subsection*{3.f Tabla de símbolos y tabla de reubicación}

\textbf{Símbolos de \texttt{main.o}:}

\begin{lstlisting}[style=asmstyle]
0000000000000000 l     F __TEXT,__text       ltmp0
0000000000000040 l     O __TEXT,__cstring    l_.str
0000000000000040 l     O __TEXT,__cstring    ltmp1
0000000000000048 l     O __LD,__compact_unwind ltmp2
0000000000000000 g     F __TEXT,__text       _main
0000000000000000         *UND*                _hello
0000000000000000         *UND*                _printf
\end{lstlisting}

\textbf{Lectura:}
\begin{itemize}[leftmargin=*]
    \item Las columnas son: dirección/offset, tipo (\texttt{l} local,
        \texttt{g} global), clase (\texttt{F} function, \texttt{O}
        object/data), sección y nombre.
    \item \texttt{\_main} está definido (\texttt{g} global, en
        \texttt{\_\_text}).
    \item \texttt{\_hello} y \texttt{\_printf} aparecen como
        \texttt{*UND*}: el módulo los \emph{usa} pero no los
        \emph{define}. El linker debe encontrarlos en otro módulo.
\end{itemize}

\textbf{Reubicaciones de \texttt{main.o}:}

\begin{lstlisting}[style=asmstyle]
RELOCATION RECORDS FOR [__text]:
OFFSET           TYPE                     VALUE
000000000000002c ARM64_RELOC_BRANCH26     _printf
0000000000000028 ARM64_RELOC_PAGEOFF12    l_.str
0000000000000024 ARM64_RELOC_PAGE21       l_.str
0000000000000018 ARM64_RELOC_BRANCH26     _hello
\end{lstlisting}

\textbf{Cómo leerlo.} Cada fila es una entrada
$reloc(instr\_address, instr\_type) \to symbol$:

\begin{itemize}[leftmargin=*, itemsep=0pt]
    \item En el offset \texttt{0x18} hay un \texttt{bl} (\emph{branch
        and link}) cuyo destino debe completarse con la dirección de
        \texttt{\_hello}.
    \item En los offsets \texttt{0x24/0x28} hay un par
        \texttt{adrp/add} que arma la dirección del literal
        \texttt{l\_.str}.
    \item En \texttt{0x2c} otro \texttt{bl}, esta vez hacia
        \texttt{\_printf}.
\end{itemize}

\textbf{Reubicaciones de \texttt{hello.o}:}

\begin{lstlisting}[style=asmstyle]
RELOCATION RECORDS FOR [__text]:
OFFSET           TYPE                     VALUE
0000000000000004 ARM64_RELOC_PAGEOFF12    l_.str
0000000000000000 ARM64_RELOC_PAGE21       l_.str
\end{lstlisting}

\textbf{Conclusión.} \texttt{hello.o} no tiene símbolos \texttt{*UND*}
(no llama a nadie); solo tiene dos entradas de reubicación para armar
la dirección de su propio literal local \texttt{l\_.str}.

\subsection*{3.g Las instrucciones \texttt{call} en \texttt{main()}}

\textbf{Procedimiento:}

\begin{lstlisting}
objdump -d main.o
\end{lstlisting}

\textbf{Resultado (\texttt{\_main}):}

\begin{lstlisting}[style=asmstyle]
0000000000000000 <ltmp0>:
   0: sub   sp, sp, #0x20
   4: stp   x29, x30, [sp, #0x10]
   8: add   x29, sp, #0x10
   c: mov   w8, #0x0
  10: str   w8, [sp, #0x8]
  14: stur  wzr, [x29, #-0x4]
  18: bl    0x18 <ltmp0+0x18>     ; <-- call a _hello (sin resolver)
  1c: mov   x8, sp
  20: str   x0, [x8]
  24: adrp  x0, 0x0 <ltmp0>
  28: add   x0, x0, #0x0
  2c: bl    0x2c <ltmp0+0x2c>     ; <-- call a _printf (sin resolver)
  30: ldr   w0, [sp, #0x8]
  ...
  3c: ret
\end{lstlisting}

\textbf{Lo importante.} Los dos \texttt{bl} apuntan ``a sí mismos'':
sus operandos son \texttt{0}. Esto coincide exactamente con lo que
describen las notas del curso: en el archivo objeto los \texttt{call}
a símbolos externos quedan sin resolver y la tabla de reubicación
guarda la información necesaria para completarlos.

\textbf{Conclusión.} \textbf{Los operandos no están resueltos en el
\texttt{.o}}: los completará el linker. Esto es exactamente lo que
veremos pasar en el ejercicio 5.

\subsection*{3.h Cómo retornan los valores las funciones}

Por la convención \textbf{AAPCS64} de ARM64:

\begin{itemize}[leftmargin=*]
    \item El valor de retorno entero o puntero se entrega en
        \texttt{x0} (o \texttt{w0} para 32 bits).
    \item La instrucción \texttt{ret} salta a la dirección guardada en
        \texttt{x30} (\emph{Link Register}), donde \texttt{bl} dejó la
        dirección de retorno.
\end{itemize}

En el código, \texttt{hello()} arma la dirección de la cadena en
\texttt{x0} y hace \texttt{ret}; \texttt{main()} guarda el puntero
recibido como argumento para \texttt{printf} y devuelve el
\texttt{int 0} en \texttt{w0} antes del \texttt{ret}. En x86-64 la
idea es la misma, pero el registro de retorno sería \texttt{rax}.

% =====================================================================
\section{Ejercicio 4 --- Generar el ejecutable \texttt{myprog}}
% =====================================================================

\textbf{Enunciado.} Combinar \texttt{main.o} y \texttt{hello.o} en un
único ejecutable y entender qué hizo el comando.

\textbf{Procedimiento:}

\begin{lstlisting}
gcc main.o hello.o -o myprog
./myprog
# Hello world
\end{lstlisting}

\subsection*{¿Se recompilaron los programas?}

\textbf{Respuesta: no.} Cuando los argumentos de \texttt{gcc} son
archivos objeto, \texttt{gcc} actúa solo como \emph{driver} del linker.
Verificación con \texttt{-v}:

\begin{lstlisting}
gcc -v main.o hello.o -o myprog_from_objects
\end{lstlisting}

\textbf{Resultado (parte relevante):}

\begin{lstlisting}[style=asmstyle]
Apple clang version 17.0.0 ...
".../usr/bin/ld" ... -o myprog_from_objects ... main.o hello.o -lSystem ...
\end{lstlisting}

No aparecen invocaciones a \texttt{cc1} o \texttt{as}: el código C
ya no se toca. Si los argumentos hubieran sido \texttt{main.c} y
\texttt{hello.c}, la traza mostraría las cuatro etapas
(preprocesador, compilador, ensamblador y linker) corriendo para cada
\texttt{.c}.

\subsection*{Qué hace el linker en concreto}

Como se discutió en la introducción, el linking tiene dos partes:

\begin{enumerate}[leftmargin=*]
    \item \textbf{Concatenar secciones y recalcular direcciones.}
        La sección \texttt{\_\_text} final del ejecutable es la
        concatenación de las \texttt{\_\_text} de \texttt{main.o} y
        \texttt{hello.o}, con direcciones absolutas únicas.
    \item \textbf{Resolver referencias externas} usando las tablas de
        símbolos y reubicación.
\end{enumerate}

En el caso concreto de \texttt{myprog}:

\begin{itemize}[leftmargin=*, itemsep=0pt]
    \item \texttt{\_main} (de \texttt{main.o}) queda en
        \texttt{0x100000460}.
    \item \texttt{\_hello} (de \texttt{hello.o}) queda en
        \texttt{0x1000004a0}, inmediatamente después.
    \item El \texttt{bl \_hello} de \texttt{main} se completa con la
        dirección real: \textbf{linking estático}.
    \item El \texttt{bl \_printf} se enlaza a un \emph{symbol stub}
        dentro del propio ejecutable que se resolverá dinámicamente
        en ejecución (ver ejercicio 5).
\end{itemize}

% =====================================================================
\section{Ejercicio 5 --- Análisis assembly de \texttt{myprog}}
% =====================================================================

\textbf{Enunciado.} Inspeccionar el ejecutable resultante y verificar
qué hizo el linker con cada referencia.

\subsection*{5.a Direcciones de \texttt{main()} y \texttt{hello()}}

\textbf{Procedimiento:}

\begin{lstlisting}
nm -n myprog
\end{lstlisting}

\textbf{Resultado:}

\begin{lstlisting}[style=outstyle]
                 U _printf
0000000100000000 T __mh_execute_header
0000000100000460 T _main
00000001000004a0 T _hello
\end{lstlisting}

\textbf{Lectura.} La columna del medio es la clase del símbolo:
\texttt{T} = texto/código definido, \texttt{U} = \emph{undefined}.

\textbf{Conclusión.} Ambos símbolos están definidos en el segmento
\texttt{\_\_TEXT}, uno detrás del otro. El linker hizo lo prometido:
tomó las dos secciones \texttt{\_\_text} (que arrancaban en offset
\texttt{0}) y las concatenó asignándoles direcciones absolutas
únicas. \texttt{\_printf} sigue siendo \texttt{U} porque su dirección
la resuelve \texttt{dyld} en tiempo de ejecución.

\subsection*{5.b ¿La llamada \texttt{main} \texttt{->} \texttt{hello} fue resuelta?}

\textbf{Respuesta: sí, estáticamente.}

\textbf{Procedimiento:}

\begin{lstlisting}
objdump -d myprog
\end{lstlisting}

\textbf{Resultado (extracto):}

\begin{lstlisting}[style=asmstyle]
_main:
100000460:  sub   sp, sp, #0x20
...
100000478:  bl    0x1000004a0 <_hello>     ; <-- llamada resuelta
10000047c:  mov   x8, sp
100000480:  str   x0, [x8]
100000484:  adrp  x0, 0x100000000
100000488:  add   x0, x0, #0x4b8           ; "%s\n"
10000048c:  bl    0x1000004ac              ; symbol stub for _printf
...
10000049c:  ret

_hello:
1000004a0:  adrp  x0, 0x100000000
1000004a4:  add   x0, x0, #0x4bc           ; "Hello world"
1000004a8:  ret
\end{lstlisting}

\textbf{Conclusión.} El \texttt{bl} en \texttt{0x100000478} ahora
codifica el desplazamiento real al inicio de \texttt{\_hello}. La
invocación quedó \emph{resuelta} porque ambos símbolos están definidos
en el mismo ejecutable.

\subsection*{5.c ¿Cómo se resolvió la invocación a \texttt{printf()}?}

\textbf{Análisis.} \texttt{printf} vive en \texttt{libSystem.B.dylib}
(la libc de macOS, equivalente a \texttt{libc.so.6} en Linux), una
biblioteca \emph{dinámica}. El linker estático no puede poner una
dirección absoluta porque el cargador puede ubicar la libc en
direcciones distintas en cada ejecución (ASLR: \emph{Address Space
Layout Randomization}).

La técnica usada es la del \textbf{symbol stub} con \emph{lazy
binding} descrita en las notas:

\begin{enumerate}[leftmargin=*]
    \item \texttt{ld} genera un stub \texttt{\_printf\$stub} en la
        sección \texttt{\_\_TEXT,\_\_stubs}.
    \item \texttt{main} hace \texttt{bl 0x1000004ac} a ese stub
        (esa dirección \emph{sí} está en el mismo ejecutable, así que
        se puede resolver estáticamente).
    \item El stub salta indirectamente a través de la
        \emph{Global Offset Table} (GOT). En Mach-O la sección se llama
        \texttt{\_\_DATA\_CONST,\_\_got}; en ELF se llama
        \texttt{.got.plt}.
    \item En la primera invocación, la GOT apunta al \emph{dynamic
        linker} (\texttt{dyld}), que resuelve la dirección real de
        \texttt{printf} en \texttt{libSystem.B.dylib} y la
        \emph{reescribe} en la GOT.
    \item Las invocaciones siguientes ya encuentran la dirección
        verdadera y saltan directamente, sin volver a llamar a
        \texttt{dyld}.
\end{enumerate}

\textbf{Verificación.}

\begin{lstlisting}
otool -Iv myprog
otool -L  myprog
\end{lstlisting}

\begin{lstlisting}[style=outstyle]
Indirect symbols for (__TEXT,__stubs) 1 entries
address            index name
0x00000001000004ac     3 _printf
Indirect symbols for (__DATA_CONST,__got) 1 entries
address            index name
0x0000000100004000     3 _printf

myprog:
    /usr/lib/libSystem.B.dylib (compatibility version 1.0.0, current version 1356.0.0)
\end{lstlisting}

\textbf{Conclusión: las dos resoluciones, lado a lado.}

\begin{center}
\small
\begin{tabular}{l l l}
\toprule
\textbf{Llamada} & \textbf{Resuelve} & \textbf{Cuándo} \\
\midrule
\texttt{main} \texttt{->} \texttt{hello}  & Linker estático (\texttt{ld})    & En tiempo de \emph{linking} \\
\texttt{main} \texttt{->} \texttt{printf} & Linker dinámico (\texttt{dyld}) & En tiempo de \emph{carga} / 1.\textsuperscript{a} llamada \\
\bottomrule
\end{tabular}
\end{center}

\begin{idea}[En Linux es lo mismo, con otros nombres]
El mismo mecanismo en ELF/Linux se llama \emph{PLT}
(Procedure Linkage Table) en lugar de \emph{symbol stub}, y la GOT
también se llama \emph{GOT}. El \emph{dynamic linker} es
\texttt{/lib64/ld-linux-x86-64.so.2} o equivalente. Las herramientas
correspondientes son \texttt{ldd} (en vez de \texttt{otool -L}) y
\texttt{readelf -d} (en vez de \texttt{otool -l}).
\end{idea}

% =====================================================================
\section{Ejercicio 6 --- Variables locales: el stack}
% =====================================================================

\textbf{Enunciado.} Demostrar experimentalmente que las variables
locales de una función residen en el \emph{stack} del proceso, y no
en \texttt{.data}/\texttt{.bss}/heap.

\textbf{Análisis previo.} El espacio de direcciones de un proceso ya
está organizado como muestra el diagrama de la introducción. Cada
invocación a una función crea un \emph{stack frame} con sus
parámetros, variables locales, dirección de retorno y registros
guardados. El compilador asigna las variables locales como
desplazamientos relativos al puntero de marco
(\texttt{x29}/\texttt{fp} en arm64), y el stack \textbf{crece hacia
direcciones bajas}: cada \texttt{call} empuja un nuevo frame hacia
abajo.

\textbf{Idea del experimento.}

\begin{enumerate}[leftmargin=*]
    \item Imprimir las direcciones de varias variables locales
        declaradas en distintas funciones.
    \item Obtener \emph{base} (dirección baja) y \emph{tope}
        (dirección alta) del stack del hilo principal con las
        funciones \texttt{pthread\_get\_stackaddr\_np} y
        \texttt{pthread\_get\_stacksize\_np}.
    \item Verificar que cada dirección de variable local cae dentro
        del rango \texttt{[base, tope)}.
\end{enumerate}

Para no romper el \texttt{main.c} usado en los demás ejercicios, el
experimento se hace en una variante: \texttt{main\_stack.c}.

\textbf{Código (extracto):}

\begin{lstlisting}[style=cstyle]
/* main_stack.c (fragmento) */
static void print_inner_locals(uintptr_t stack_base, uintptr_t stack_top)
{
    int inner_a = 10;
    int inner_b = 20;
    char inner_buf[16] = "inner";
    report_local("&inner_a",  &inner_a,  stack_base, stack_top);
    report_local("&inner_b",  &inner_b,  stack_base, stack_top);
    report_local("inner_buf", inner_buf, stack_base, stack_top);
}

static void print_stack_layout(void)
{
    int    local_a = 1;
    int    local_b = 2;
    double local_c = 3.14;
    char   local_d = 'x';
    char   buffer[32] = "stack-local";
    pthread_t self = pthread_self();
    void   *stack_top  = pthread_get_stackaddr_np(self);
    size_t  stack_size = pthread_get_stacksize_np(self);
    void   *stack_base = (void *)((uintptr_t)stack_top - stack_size);
    /* imprimir base, tope, tamano y direcciones de cada local */
    print_inner_locals((uintptr_t)stack_base, (uintptr_t)stack_top);
}
\end{lstlisting}

\textbf{Procedimiento:}

\begin{lstlisting}
gcc -Wall -Wextra -pedantic main_stack.c hello.c -o myprog_stack
./myprog_stack
\end{lstlisting}

\textbf{Resultado.} Las direcciones absolutas cambian entre
ejecuciones por ASLR, pero la \emph{estructura relativa} se mantiene:

\begin{lstlisting}[style=outstyle]
Hello world
Stack del hilo principal:
  base (dir baja): 0x16ef98000
  tope (dir alta): 0x16f794000
  tamano: 8372224 bytes (7.98 MiB)

Locales en main/print_stack_layout:
  &local_a       0x16f7924fc  [STACK]
  &local_b       0x16f7924f8  [STACK]
  &local_c       0x16f7924f0  [STACK]
  &local_d       0x16f7924ef  [STACK]
  buffer         0x16f792500  [STACK]

Locales en funcion secundaria:
  &inner_a       0x16f79243c  [STACK]
  &inner_b       0x16f792438  [STACK]
  inner_buf      0x16f792450  [STACK]
\end{lstlisting}

\textbf{Análisis de la salida.}

\begin{itemize}[leftmargin=*]
    \item Todas las direcciones impresas caen dentro del rango
        \texttt{[stack\_base, stack\_top)}: las variables locales
        viven, efectivamente, en el \textbf{stack}.
    \item El stack del hilo principal mide $\approx 8$\,MiB
        (\texttt{8372224} bytes), valor por defecto en macOS y similar
        al de Linux (\texttt{ulimit -s} típicamente da 8192 KiB).
    \item Las locales de \texttt{print\_inner\_locals} están en
        direcciones \emph{más bajas} (\texttt{0x16f79243c}) que las de
        \texttt{print\_stack\_layout} (\texttt{0x16f7924fc}): cada
        invocación empujó un nuevo frame ``hacia abajo''.
    \item El alineamiento se respeta: \texttt{local\_c} (un
        \texttt{double}) cae en una dirección múltiplo de 8.
\end{itemize}

\begin{ojo}[Una variable local con \texttt{static} no va al stack]
Si declaráramos \texttt{static int local\_a = 1;} dentro de una
función, \texttt{local\_a} no estaría en el stack: viviría en
\texttt{.data} (porque tiene valor inicial) como una global de visibilidad
restringida. Lo mismo para variables globales fuera de cualquier
función. Es importante distinguir \emph{duración de almacenamiento}
(automática, estática, dinámica) de \emph{scope}.
\end{ojo}

\textbf{Respuesta concreta.} Las variables locales se asignan en el
\emph{stack} del proceso, una región del espacio de direcciones que
crece hacia direcciones bajas y se organiza en \emph{stack frames}
(uno por invocación de función). ASLR aleatoriza las direcciones
absolutas en cada ejecución, pero la estructura relativa (frames
internos por debajo de los externos) permanece.

% =====================================================================
\section{Ejercicio 7 --- Biblioteca estática \texttt{libhello.a}}
% =====================================================================

\textbf{Enunciado.} Empaquetar \texttt{hello.o} y un segundo módulo
\texttt{hello2.o} en una biblioteca estática, y enlazar un programa
contra ella.

\textbf{Análisis.} Una biblioteca \textbf{estática} es un \emph{archive}
(formato \texttt{.a}) que agrupa varios \texttt{.o} junto con un
\emph{índice de símbolos}. El linker, ante una referencia sin
resolver, busca en el índice del \texttt{.a} y, si lo encuentra,
\textbf{copia el \texttt{.o} correspondiente al ejecutable}.

Para no romper \texttt{main.c}, este ejercicio usa una variante
(\texttt{main\_lib.c}) y un segundo módulo (\texttt{hello2.c}):

\begin{lstlisting}[style=cstyle]
/* hello2.c */
#define hi2 "Hello from hello2"
char *hello2(void) { return hi2; }
\end{lstlisting}

\begin{lstlisting}[style=cstyle]
/* main_lib.c */
#include <stdio.h>
extern char *hello(void);
extern char *hello2(void);
int main(void) {
    printf("%s\n", hello());
    printf("%s\n", hello2());
    return 0;
}
\end{lstlisting}

\subsection*{7.a Crear la biblioteca estática}

\textbf{Procedimiento:}

\begin{lstlisting}
gcc -Wall -Wextra -pedantic -c hello.c  -o hello.o
gcc -Wall -Wextra -pedantic -c hello2.c -o hello2.o
ar rcs libhello.a hello.o hello2.o
\end{lstlisting}

\textbf{Análisis de las opciones de \texttt{ar}.}
\begin{itemize}[leftmargin=*, itemsep=0pt]
    \item \texttt{r} \emph{reemplaza} miembros existentes (o los
        agrega si no existían).
    \item \texttt{c} \emph{crea} el archivo si no existe (silencia el
        warning).
    \item \texttt{s} genera o actualiza el \emph{índice de símbolos}.
        Es lo que permite al linker responder rápidamente ``¿qué
        objeto define a \texttt{hello}?''.
\end{itemize}

\subsection*{7.b Listar los módulos de la biblioteca}

\textbf{Procedimiento:}

\begin{lstlisting}
ar -t libhello.a
\end{lstlisting}

\textbf{Resultado:}

\begin{lstlisting}[style=outstyle]
__.SYMDEF SORTED
hello.o
hello2.o
\end{lstlisting}

\textbf{Conclusión.} La primera entrada (\texttt{\_\_.SYMDEF SORTED})
es el índice de símbolos interno generado por \texttt{ar -s}; los
módulos reales son \texttt{hello.o} y \texttt{hello2.o}. En ELF/Linux
el índice equivalente se llama \texttt{/SYMDEF} o
\texttt{\_\_index} según implementación.

\subsection*{7.c Compilar y linkear contra la biblioteca}

\textbf{Procedimiento:}

\begin{lstlisting}
gcc -Wall -Wextra -pedantic main_lib.c -L. -lhello -o myprog_lib
./myprog_lib
\end{lstlisting}

\textbf{Resultado:}

\begin{lstlisting}[style=outstyle]
Hello world
Hello from hello2
\end{lstlisting}

\textbf{Detalle del comando.}
\begin{itemize}[leftmargin=*, itemsep=0pt]
    \item \texttt{-L.} agrega el directorio actual al \emph{path} de
        búsqueda de bibliotecas.
    \item \texttt{-lhello} pide buscar una biblioteca llamada
        \texttt{libhello}. El prefijo \texttt{lib} y el sufijo
        (\texttt{.a} / \texttt{.dylib} / \texttt{.so}) los pone el
        propio linker.
\end{itemize}

\begin{ojo}[Si conviven \texttt{.a} y \texttt{.dylib}, gana la dinámica]
Si en el mismo directorio existieran \texttt{libhello.a} y
\texttt{libhello.dylib}, \texttt{ld} prefiere por defecto la dinámica.
Es lo que ocurre en este directorio: al compilar con \texttt{-lhello}
el linker elige \texttt{libhello.dylib}, que solo exporta
\texttt{hello}, y la compilación falla por \texttt{hello2}
no encontrado. Para forzar la estática hay dos opciones: pasar
\texttt{libhello.a} directamente en la línea de comandos
(\texttt{gcc main\_lib.c libhello.a}), o usar la flag
\texttt{-Wl,-search\_paths\_first} con \texttt{libhello.a} explícita.
En Linux la flag equivalente es \texttt{-Wl,-Bstatic ... -Wl,-Bdynamic}.
\end{ojo}

\textbf{Conclusión.} El linker, al ver que \texttt{main\_lib.o}
referencia \texttt{\_hello} y \texttt{\_hello2}, consulta el índice
de \texttt{libhello.a}, encuentra que provienen de \texttt{hello.o} y
\texttt{hello2.o}, \textbf{extrae esos miembros y los integra} en el
ejecutable. El resultado es un binario \emph{autocontenido}: una vez
linkeado no necesita la biblioteca para correr.

% =====================================================================
\section{Ejercicio 8 --- Biblioteca dinámica y \texttt{dlopen}}
% =====================================================================

\textbf{Enunciado.} Generar una biblioteca dinámica con la función
\texttt{hello}, escribir un programa que la cargue \emph{en tiempo de
ejecución} con \texttt{dlopen}/\texttt{dlsym}, y verificar que el
ejecutable \emph{no} contiene el código de \texttt{hello}.

\textbf{Análisis.} Las bibliotecas \textbf{dinámicas}
(\texttt{.so}/\texttt{.dylib}/\texttt{.dll}) son archivos que el SO
carga y enlaza bajo demanda. El ejecutable no incluye su código: si
varios procesos las usan, una sola copia en memoria los sirve a todos.

Existen dos formas de usarlas:

\begin{itemize}[leftmargin=*]
    \item \textbf{Linking dinámico implícito} (lo que vimos con
        \texttt{printf} en el ejercicio 5): el linker estático genera
        \emph{stubs} y entradas en la GOT al momento de enlazar el
        ejecutable. El \emph{dynamic linker} resuelve los símbolos al
        cargar el programa.
    \item \textbf{Carga dinámica explícita} (este ejercicio): el
        programa decide en tiempo de ejecución qué biblioteca cargar,
        usando la API POSIX \texttt{<dlfcn.h>}
        (\texttt{dlopen}/\texttt{dlsym}/\texttt{dlclose}). Es la base
        de los \emph{plugins} y de cualquier mecanismo de
        extensibilidad en ejecución.
\end{itemize}

\subsection*{8.a Generar la biblioteca dinámica}

\textbf{Procedimiento:}

\begin{lstlisting}
gcc -Wall -Wextra -pedantic -fPIC -shared hello.c -o libhello.dylib
file libhello.dylib
\end{lstlisting}

\textbf{Resultado:}

\begin{lstlisting}[style=outstyle]
libhello.dylib: Mach-O 64-bit dynamically linked shared library arm64
\end{lstlisting}

\textbf{Por qué cada flag importa.}
\begin{itemize}[leftmargin=*, itemsep=0pt]
    \item \texttt{-fPIC} (\emph{Position Independent Code}) es
        obligatorio: la biblioteca puede cargarse en distintas
        direcciones en cada proceso (ASLR), así que el código no
        puede asumir una dirección de carga fija. Toda referencia a
        un dato global pasa por la GOT.
    \item \texttt{-shared} le indica a \texttt{gcc} que la salida es
        una biblioteca compartida, no un ejecutable.
\end{itemize}

\textbf{Verificación de los símbolos exportados:}

\begin{lstlisting}
nm libhello.dylib
\end{lstlisting}

\begin{lstlisting}[style=outstyle]
0000000000000320 T _hello
\end{lstlisting}

La biblioteca exporta exactamente un símbolo: \texttt{\_hello}.

\subsection*{8.b/8.c Carga dinámica con \texttt{dlopen}}

\textbf{Código:}

\begin{lstlisting}[style=cstyle]
/* main_dlopen.c */
#include <stdio.h>
#include <stdlib.h>
#include <dlfcn.h>

typedef char *(*hello_fn_t)(void);

int main(void)
{
    const char *lib_path = "./libhello.dylib";
    void *handle = dlopen(lib_path, RTLD_NOW);
    if (handle == NULL) {
        fprintf(stderr, "dlopen: %s\n", dlerror());
        return 1;
    }
    dlerror();   /* limpia el estado de error previo */

    hello_fn_t hello_fn = NULL;
    *(void **)(&hello_fn) = dlsym(handle, "hello");
    const char *err = dlerror();
    if (err != NULL) {
        fprintf(stderr, "dlsym: %s\n", err);
        dlclose(handle);
        return 1;
    }

    printf("%s\n", hello_fn());

    dlclose(handle);
    return 0;
}
\end{lstlisting}

\textbf{Decisiones de diseño.}
\begin{itemize}[leftmargin=*]
    \item \texttt{RTLD\_NOW} fuerza la resolución de todos los
        símbolos al cargar la biblioteca. \texttt{RTLD\_LAZY} los
        resolvería bajo demanda; para una biblioteca pequeña es
        indistinto.
    \item \texttt{*(void **)(\&hello\_fn) = dlsym(...)} es el modismo
        recomendado por POSIX para evitar el \emph{warning} sobre
        conversión de \texttt{void *} a puntero a función (en C
        estándar esa conversión no está garantizada, aunque en POSIX
        sí lo está).
    \item Se llama a \texttt{dlerror()} antes y después de
        \texttt{dlsym} para no confundir un símbolo con valor
        legítimamente \texttt{NULL} con un símbolo \emph{no
        encontrado}.
\end{itemize}

\textbf{Compilación y ejecución}
(en macOS la libc provee \texttt{dlopen} de fábrica;
en Linux haría falta agregar \texttt{-ldl}):

\begin{lstlisting}
gcc -Wall -Wextra -pedantic main_dlopen.c -o myprog_dlopen
./myprog_dlopen
\end{lstlisting}

\begin{lstlisting}[style=outstyle]
Hello world
\end{lstlisting}

\subsection*{8.c El código de \texttt{hello} no está en el ejecutable}

\textbf{Procedimiento:}

\begin{lstlisting}
otool -L myprog_dlopen
nm     myprog_dlopen | grep hello
\end{lstlisting}

\textbf{Resultado:}

\begin{lstlisting}[style=outstyle]
myprog_dlopen:
    /usr/lib/libSystem.B.dylib (compatibility version 1.0.0, current version 1356.0.0)

(la busqueda de "hello" en nm no devuelve nada)
\end{lstlisting}

\textbf{Conclusión.} La única dependencia es
\texttt{libSystem.B.dylib} (que provee \texttt{printf},
\texttt{dlopen}, etc.). \textbf{No hay referencia a
\texttt{libhello.dylib}} porque el ejecutable nunca fue enlazado
contra ella; la cargará el propio programa con \texttt{dlopen}. Y
\texttt{nm} no encuentra el símbolo \texttt{\_hello} porque ese código
no está en el binario: vive en la biblioteca dinámica y se carga en
ejecución.

\subsection*{8.d Sección de enlazado dinámico}

En Linux se usaría \texttt{readelf -d myprog\_dlopen} para ver la
sección \texttt{.dynamic}. En macOS el equivalente es \texttt{otool}:

\begin{lstlisting}
otool -L libhello.dylib
\end{lstlisting}

\begin{lstlisting}[style=outstyle]
libhello.dylib:
    libhello.dylib (compatibility version 0.0.0, current version 0.0.0)
    /usr/lib/libSystem.B.dylib ...
\end{lstlisting}

La biblioteca declara su propio nombre como \emph{install name} y a
su vez depende de la libc (\texttt{libSystem.B.dylib}).

\subsection*{Conexión con la teoría}

\begin{idea}[Implícito vs explícito]
\begin{itemize}[leftmargin=*, itemsep=0pt]
    \item \textbf{Linking dinámico implícito} (ejercicios 4 y 5): el
        linker estático genera stubs y entradas en la GOT, y el
        \emph{dynamic linker} resuelve los símbolos en tiempo de
        carga sin que el programa tenga que hacer nada.
    \item \textbf{Carga dinámica explícita} (este ejercicio): el
        propio programa decide qué biblioteca cargar y cuándo,
        mediante \texttt{dlopen}/\texttt{dlsym}. Es lo que usan los
        sistemas de \emph{plugins}, los lenguajes interpretados
        cuando cargan extensiones, y las extensiones del kernel
        (módulos).
\end{itemize}
\end{idea}

% =====================================================================
\section{Ejercicio 9 --- Automatización con \texttt{Makefile}}
% =====================================================================

\textbf{Enunciado.} Empaquetar la construcción de \texttt{myprog} en
un \texttt{Makefile} que permita compilación incremental, y verificar
que \texttt{make} solo reconstruye lo que cambió.

\textbf{Análisis.} En un proyecto real, los pasos del toolchain no se
invocan a mano: se automatizan con un \emph{build system}. El más
clásico en UNIX es \textbf{GNU make}, basado en un \texttt{Makefile}
que describe \emph{objetivos} (\emph{targets}), sus
\emph{dependencias} y los \emph{comandos} para construirlos. La forma
básica de una regla es

\begin{lstlisting}[style=asmstyle]
target: dependencies
<TAB>command
\end{lstlisting}

\texttt{make} construye un grafo de dependencias y, comparando los
\emph{timestamps} de cada archivo, decide qué reglas disparar. Si el
objetivo es más nuevo que todas sus dependencias, no se reconstruye:
de ahí la \textbf{compilación incremental}.

Las \textbf{variables automáticas} clave son:
\begin{itemize}[leftmargin=*, itemsep=0pt]
    \item \texttt{\$@}: nombre del target.
    \item \texttt{\$\^{}}: todas las dependencias.
    \item \texttt{\$<}: la primera dependencia.
\end{itemize}

Y las \textbf{reglas de patrón} (\texttt{\%.o: \%.c}) evitan repetir
la misma acción para cada archivo fuente.

\subsection*{9.1 Makefile usado}

Los originales \texttt{main.c} y \texttt{hello.c} viven en el directorio
del práctico. Para mantener la convención de los demás ejercicios, el
\texttt{Makefile} se ubica dentro de \texttt{ejercicio-9/} y referencia
los fuentes vía la variable especial \texttt{VPATH}.

\begin{lstlisting}[style=asmstyle]
# VPATH le dice a make donde buscar los archivos fuente cuando no
# estan en el directorio actual.
VPATH = ..

CC     = gcc
CFLAGS = -Wall -Wextra -pedantic

# Objetivo por defecto: myprog
myprog: main.o hello.o
	$(CC) $(CFLAGS) -o $@ $^

main.o: main.c
	$(CC) $(CFLAGS) -c $< -o $@

hello.o: hello.c
	$(CC) $(CFLAGS) -c $< -o $@

.PHONY: clean
clean:
	rm -f main.o hello.o myprog
\end{lstlisting}

\textbf{Decisiones de diseño.}
\begin{itemize}[leftmargin=*]
    \item \texttt{VPATH = ..} es la forma idiomática de separar
        fuentes y artefactos de compilación: \texttt{make} encuentra
        \texttt{main.c}/\texttt{hello.c} un nivel arriba y produce
        los \texttt{.o} y el ejecutable en el directorio actual.
    \item \texttt{CC} y \texttt{CFLAGS} centralizan las opciones del
        compilador.
    \item \texttt{\$@}, \texttt{\$\^{}}, \texttt{\$<} hacen las reglas
        concisas y reutilizables.
    \item \texttt{.PHONY: clean} declara \texttt{clean} como objetivo
        ``ficticio'': si existiera un archivo llamado \texttt{clean},
        sin esta declaración \texttt{make} lo consideraría ``ya
        construido'' y no ejecutaría la regla.
\end{itemize}

\subsection*{9.2 Grafo de dependencias}

\begin{center}
\begin{tikzpicture}[
    node distance=0.9cm,
    every node/.style={draw, rounded corners, font=\ttfamily\small,
                       minimum width=1.6cm, align=center},
    arr/.style={-{Latex[length=2mm]}, thick}
]
\node (myprog) {myprog};
\node (mainO)  [below left=of myprog]  {main.o};
\node (helloO) [below right=of myprog] {hello.o};
\node (mainC)  [below=of mainO]  {main.c};
\node (helloC) [below=of helloO] {hello.c};
\draw[arr] (mainO)  -- (myprog);
\draw[arr] (helloO) -- (myprog);
\draw[arr] (mainC)  -- (mainO);
\draw[arr] (helloC) -- (helloO);
\end{tikzpicture}
\end{center}

\texttt{make myprog} requiere \texttt{main.o} y \texttt{hello.o};
cada \texttt{.o} requiere su \texttt{.c}. \texttt{make} recorre el
grafo en post-orden: primero compila los \texttt{.o} que falten o
estén desactualizados y solo entonces invoca al linker.

\subsection*{9.3 Ejecución desde cero}

\begin{lstlisting}[style=outstyle]
$ make clean
rm -f main.o hello.o myprog

$ make
gcc -Wall -Wextra -pedantic -c ../main.c -o main.o
gcc -Wall -Wextra -pedantic -c ../hello.c -o hello.o
gcc -Wall -Wextra -pedantic -o myprog main.o hello.o

$ ./myprog
Hello world
\end{lstlisting}

\subsection*{9.4 Compilación incremental}

Si se vuelve a invocar \texttt{make} sin tocar nada, el ejecutable ya
está al día y \texttt{make} no ejecuta ningún comando:

\begin{lstlisting}[style=outstyle]
$ make
make: `myprog' is up to date.
\end{lstlisting}

Si se modifica solo un archivo fuente, \texttt{make} rehace
\textbf{únicamente lo necesario}:

\begin{lstlisting}[style=outstyle]
$ touch ../hello.c
$ make
gcc -Wall -Wextra -pedantic -c ../hello.c -o hello.o
gcc -Wall -Wextra -pedantic -o myprog main.o hello.o
\end{lstlisting}

\texttt{main.o} no se recompiló porque \texttt{main.c} no cambió. En
proyectos grandes, recompilar todo desde cero es inviable;
\texttt{make} escala porque solo toca lo que tiene que tocar.

\subsection*{9.5 Versión compacta con regla de patrón}

La misma funcionalidad puede escribirse aún más breve usando una regla
de patrón \texttt{\%.o : \%.c}, que define cómo construir
\emph{cualquier} \texttt{.o} a partir de un \texttt{.c}:

\begin{lstlisting}[style=asmstyle]
VPATH  = ..
CC     = gcc
CFLAGS = -Wall -Wextra -pedantic

myprog: main.o hello.o
	$(CC) $(CFLAGS) -o $@ $^

%.o: %.c
	$(CC) $(CFLAGS) -c $< -o $@

.PHONY: clean
clean:
	rm -f main.o hello.o myprog
\end{lstlisting}

La versión con reglas explícitas (la que está activa) es más
didáctica porque hace visibles los tres comandos del toolchain. La
versión con patrón es la que se usa en proyectos reales porque escala
sin repetir reglas.

\textbf{Conclusión.} El \texttt{Makefile} no agrega capacidades nuevas
al proceso de compilación; \emph{automatiza} exactamente las mismas
etapas vistas en los ejercicios 1--4. Lo que aporta es
\emph{automatización con conocimiento de dependencias}, la
herramienta que hace manejables los proyectos grandes manteniendo
compilaciones rápidas y reproducibles.

% =====================================================================
\section*{Cierre}
\addcontentsline{toc}{section}{Cierre}
% =====================================================================

A lo largo del práctico se ``abrió'' el proceso que normalmente queda
oculto detrás del comando \texttt{gcc -o myprog \dots}:

\begin{itemize}[leftmargin=*]
    \item Pre-procesamiento, compilación y ensamblado generan los
        artefactos intermedios (\texttt{.i}, \texttt{.s}, \texttt{.o})
        inspeccionables con \texttt{objdump} / \texttt{otool} /
        \texttt{nm}.
    \item Los archivos objeto contienen secciones (\texttt{.text},
        \texttt{.data}, \texttt{.bss}, \texttt{.rodata}/\texttt{\_\_cstring}),
        tabla de símbolos y tabla de reubicación; las funciones
        empiezan en offset \texttt{0x0} y los \texttt{call} a símbolos
        externos quedan sin resolver.
    \item El linker concatena las secciones, recalcula direcciones y
        resuelve referencias: \emph{estáticamente} las que se conocen
        (como \texttt{hello}) y \emph{dinámicamente} las que viven en
        bibliotecas compartidas (como \texttt{printf}, vía
        stubs/PLT y GOT).
    \item Las variables locales se asignan en el \emph{stack}, una
        región del proceso que crece hacia direcciones bajas; se
        verificó experimentalmente que las direcciones impresas caen
        dentro del rango reportado por
        \texttt{pthread\_get\_stackaddr\_np}.
    \item Las bibliotecas estáticas (\texttt{libhello.a}) son
        contenedores \texttt{ar} de \texttt{.o} con índice de
        símbolos; el linker extrae solo los miembros necesarios y los
        integra al ejecutable.
    \item Las bibliotecas dinámicas
        (\texttt{libhello.dylib}/\texttt{libhello.so}) se cargan en
        tiempo de ejecución, ya sea de forma implícita (vía stubs y
        GOT) o explícita (vía \texttt{dlopen}/\texttt{dlsym}).
    \item \texttt{make} automatiza todo el flujo: un Makefile declara
        objetivos, dependencias y comandos, y \texttt{make}
        reconstruye solo lo que cambió comparando \emph{timestamps}.
\end{itemize}

Esta maquinaria es la base de los conceptos de los prácticos
siguientes: cómo el SO carga un ejecutable en memoria, cómo organiza
el espacio de direcciones del proceso y cómo distintos procesos
pueden compartir código mediante bibliotecas dinámicas.

% =====================================================================
\section*{Referencias}
\addcontentsline{toc}{section}{Referencias}
% =====================================================================

\begin{enumerate}[leftmargin=*]
    \item Marcelo Arroyo. \emph{Notas del curso de Sistemas
        Operativos}, capítulo \emph{Herramientas de desarrollo}.
        UNRC, 2026.
        \url{https://marceloarroyo.gitlab.io/cursos/SO/index.html\#/toolchain}
    \item Brian Kernighan, Dennis Ritchie. \emph{The C Programming
        Language}, 2.\textsuperscript{a} edición. Prentice Hall.
    \item John R.\ Levine. \emph{Linkers and Loaders}. Morgan
        Kaufmann, 2000.
    \item Páginas de manual: \texttt{man 1 gcc}, \texttt{man 1 ld},
        \texttt{man 1 ar}, \texttt{man 1 objdump},
        \texttt{man 1 otool}, \texttt{man 1 nm},
        \texttt{man 1 readelf}, \texttt{man 1 ldd},
        \texttt{man 3 dlopen}.
\end{enumerate}

\end{document}
